\section{Related Work}\label{sec:related}

Automated PPE detection has consolidated around the YOLO family of one-stage detectors, whose inference speed makes them the default choice for video-rate safety applications \cite{ultralytics2023}. Recent YOLOv8-based studies benchmark helmet and vest detection across public object-detection datasets \cite{ppeyolo2024}, and edge-deployed variants extend the approach to preventing falls from height on construction sites, confirming that real-time PPE monitoring is feasible on modest hardware \cite{edgeyolo2024}. Community resources such as the Roboflow Universe PPE-Detection corpus \cite{roboflowppe} supply most of the training data behind this line of work and have made baseline detectors straightforward to reproduce. The shared limitation is scope. These benchmarks are construction-centric: their class inventories revolve around hard hats and high-visibility vests, while the respiratory and eye protection that dominates compliance requirements in chemical processing---gas masks, goggles---is absent. None of them adapts the detector to the visual conditions of a specific site, where dust, harsh lighting, and uniforms that resemble PPE degrade models trained on clean public imagery.

The SH17 dataset \cite{ahmad2024sh17} broadens the label space to seventeen classes spanning PPE items and human body parts in manufacturing scenes, currently the widest public inventory for industrial PPE detection. Its breadth, however, is deliberately generic: the imagery reflects general manufacturing rather than chemical-industry environments, and the dataset carries no mechanism for the domain shift a detector encounters when moved onto a particular plant's cameras. SH17 is a strong ingredient for a site-adapted model, not a substitute for one.

A separate literature explains why automated detection is needed at all. Classical sustained-attention research established that vigilance decays within the first half hour of a continuous monitoring task \cite{parasuraman1984}, and that the decrement persists even when observers are explicitly instructed and motivated \cite{sawin1995}. CCTV-specific evidence points the same way. Velastin et al. proposed motion-based video analysis for underground stations on the premise that human operators cannot reliably watch large banks of monitors \cite{velastin2006}, and a Sandia National Laboratories assessment of continuous monitoring tasks describes attention lapses that spread across operators as sessions extend \cite{sandia2014}. The most direct evidence for industrial CCTV comes from Donald, Donald, and Thatcher, who studied trained surveillance operators in diamond-processing plants and reported that they detected roughly half of the targeted risk behaviours while generating a high rate of false alarms \cite{donald2015}. That second finding matters as much as the first. A detector that merely replaces the human eye trades missed events for a stream of false alerts, shifting the vigilance burden downstream to whoever triages the alarms rather than removing it. The false-alarm half of that result is the specific failure mode our verification stage is built against: no alert should reach safety staff before a reasoning pass has examined why the detection might be wrong.

Vision-language models are a natural candidate for that verification role. Recent reviews chart rapid progress in VLMs for object detection and segmentation \cite{vlmreview2025}, and open-weight models such as Qwen2.5-VL bring grounded visual reasoning to hardware that an industrial operator can run on premises \cite{qwen25vl}. Detect-then-reason pipelines, in which a fast detector proposes candidate events and a slower model judges them with scene context, follow directly from these capabilities. Yet the pattern remains nascent, and applications to safety monitoring---where occlusion, glare, and distance routinely produce the kind of ambiguous detection that contextual reasoning could resolve---are rare.

Each thread therefore stops short of our setting. Detection research offers fast, reproducible models but only generic, construction-centric classes and no site adaptation; SH17 widens the label space without reaching chemical-industry classes or conditions; the vigilance literature diagnoses manual monitoring as structurally unreliable and warns that false alarms recreate the problem downstream, but offers no detection-side remedy; and VLM-based verification exists in principle but not in deployed safety pipelines. No prior work, to our knowledge, joins site-adapted detection over chemical-industry PPE classes with a reasoning-based false-positive filter on retrofit CCTV infrastructure. Table~\ref{tab:relwork} summarizes the closest prior efforts and what each contributes to that combined gap.

\begin{table}[t]
\caption{Summary of Related Work Across the Three Threads}
\label{tab:relwork}
\centering
\footnotesize
\setlength{\tabcolsep}{3pt}
\begin{tabular}{p{1.9cm}p{1.6cm}p{1.9cm}p{2.5cm}}
\toprule
\textbf{Study/Source} & \textbf{Method} & \textbf{Focus} & \textbf{Key Finding} \\
\midrule
Roboflow PPE-Detection \cite{roboflowppe} & Public YOLO-format dataset & Helmet/vest detection & Widely reused baseline; construction-centric class set \\
\addlinespace
SH17 \cite{ahmad2024sh17} & Object-detection dataset & 17 PPE and body-part classes in manufacturing & Broadest public industrial PPE label space; generic conditions \\
\addlinespace
PPE-YOLOv8 comparative study \cite{ppeyolo2024} & YOLOv8 benchmarking & PPE detection on public benchmark datasets & YOLOv8 performs competitively for PPE detection \\
\addlinespace
Edge-YOLOv8 worker safety \cite{edgeyolo2024} & Improved YOLOv8 on edge devices & Fall-from-height PPE, construction sites & Real-time PPE detection feasible on edge hardware \\
\addlinespace
Donald et al. (2015) \cite{donald2015} & Field study of CCTV operators & Operator performance, diamond-processing plants & Reported low detection of target behaviours with high false-alarm rates \\
\addlinespace
Velastin et al. (2006) \cite{velastin2006} & Motion-based video analysis & Automated surveillance support & Automation proposed to offset the limits of manual multi-camera monitoring \\
\addlinespace
VLM review (2025) \cite{vlmreview2025} & Survey and evaluation & VLMs for detection and segmentation & Rapid progress; safety-monitoring applications still nascent \\
\bottomrule
\end{tabular}
\end{table}
